\documentclass{beamer}

\usetheme{Berlin}
\usecolortheme{Orchid}
\useoutertheme{miniframes}
\setbeamertemplate{navigation symbols}{}

\usepackage[utf8]{inputenc}
\usepackage[T1]{fontenc}
\usepackage{amsmath}
\usepackage{amssymb}
\usepackage{booktabs}
\usepackage{graphicx}
\graphicspath{{../images/}}

\usepackage{xcolor}
\usepackage{listings}
\usepackage{hyperref}
\usepackage{tikz}
\usetikzlibrary{arrows.meta,positioning,shapes.geometric,calc}

\lstset{
  language=Java,
  basicstyle=\ttfamily\small,
  keywordstyle=\color{blue},
  commentstyle=\color{gray},
  stringstyle=\color{teal},
  showstringspaces=false,
  columns=fullflexible,
  breaklines=true
}

% Per-slide source line (references shown on the slide itself).
\newcommand{\slidesources}[1]{%
  \vfill
  {\tiny\color{gray}\raggedright\sloppy\parbox{\linewidth}{Sources: #1}\par}%
}

% Unit-wise source strings for this subject (used by slides/notes).
% Path is relative to the lecture `latex/` folder (the compilation CWD).
% OOPS with Java (CSEG1044) — unit-wise sources
%
% These are short, student-facing reference strings intended to be used with:
%   - \slidesources{...}  (in beamer slides)
%   - \notesources{...}   (in lecture notes)
%
% Style inspiration: other subjects in My-Subjects/ use semicolon-separated sources
% and include an "accessed" date for web documentation.

\newcommand{\OOPSUnitISources}{Schildt, \textit{Java: A Beginner's Guide} (10e), McGraw-Hill (Java fundamentals + OOP basics).; Oracle Java Tutorials: Getting Started + Language Basics + Classes and Objects (accessed \today).; Java SE API docs (e.g., \texttt{String}, \texttt{Scanner}) (accessed \today).}

\newcommand{\OOPSUnitIISources}{Schildt, \textit{Java: The Complete Reference} (12e), McGraw-Hill (classes, inheritance, packages, interfaces).; Bloch, \textit{Effective Java} (3e), Addison-Wesley, 2018 (design choices; interfaces vs abstract classes).; Oracle Java Tutorials: Classes and Objects + Interfaces + Packages (accessed \today).; Java SE API docs (e.g., \texttt{Comparable}, \texttt{Runnable}) (accessed \today).}

\newcommand{\OOPSUnitIIISources}{Schildt, \textit{Java: The Complete Reference} (12e) (nested/inner classes, exceptions, threads, I/O).; Oracle Java Tutorials: Nested Classes + Exceptions + Concurrency + Basic I/O (accessed \today).; Java SE API docs (e.g., \texttt{Thread}, \texttt{AutoCloseable}) (accessed \today).}

\newcommand{\OOPSUnitIVSources}{Schildt, \textit{Java: The Complete Reference} (12e) (generics, lambdas, Swing, JDBC).; Bloch, \textit{Effective Java} (3e) (generics + lambdas).; Oracle Java Tutorials: Generics + Lambda Expressions + Swing + JDBC Basics (accessed \today).; Java SE API docs (e.g., \texttt{java.util.function}, \texttt{javax.swing}, \texttt{java.sql}) (accessed \today).}

\newcommand{\OOPSUnitVSources}{Schildt, \textit{Java: The Complete Reference} (12e) (Collections Framework + wrapper classes).; Oracle Java docs: Collections Framework overview (accessed \today).; Java SE API docs (\texttt{java.util} collections; wrapper classes in \texttt{java.lang}) (accessed \today).}

\newcommand{\OOPSUnitVISources}{Git SCM: \textit{Pro Git} (2e) (version control workflow), accessed \today.; GitHub Docs (repositories, issues, pull requests), accessed \today.; Oracle Java Tutorials: Swing + JDBC (accessed \today).; JUnit 5 User Guide (accessed \today).; Apache Maven Project: Getting Started (accessed \today).; Gradle User Manual: Getting Started (accessed \today).; UML class diagrams: UML 2.x overview/spec (accessed \today).}

\newcommand{\OOPSMiscWhyInterfacesSources}{Schildt, \textit{Java: The Complete Reference} (12e) (interfaces).; Bloch, \textit{Effective Java} (3e) (prefer interfaces where appropriate).; Oracle Java Tutorials: Interfaces (accessed \today).; Java SE API docs (e.g., \texttt{Comparable}, \texttt{Runnable}) (accessed \today).}



\title[OOPS with Java]{Object Oriented Programming (CSEG1044)}
\subtitle{Unit 05 -- Lecture 03: Wrapped Collections + `Collections` Utilities + Wrappers}
\begin{document}

\begin{frame}[fragile]
  \titlepage
  \vspace{0.5em}
  \slidesources{\OOPSUnitVSources}
\end{frame}

\begin{frame}{Quick Links}
  \centering
  \hyperlink{sec:overview}{\beamerbutton{Overview}}\hspace{0.6em}
  \hyperlink{sec:concepts}{\beamerbutton{Key Topics}}\hspace{0.6em}
  \hyperlink{sec:demo}{\beamerbutton{Demo}}\hspace{0.6em}
  \hyperlink{sec:summary}{\beamerbutton{Summary}}
  \slidesources{\OOPSUnitVSources}
\end{frame}

\begin{frame}{Agenda}
  \tableofcontents
  \slidesources{\OOPSUnitVSources}
\end{frame}

\section{Overview}
\label{sec:overview}

\begin{frame}{Lecture Snapshot}
\begin{itemize}[<+->]
  \item \textbf{Unit focus:} Collections Framework and Wrapper Classes
  \item \textbf{Course thread:} Use Wrapped Collections + `Collections` Utilities + Wrappers to manage in-memory project data more cleanly and predictably inside Campus Resource Manager.
\end{itemize}
  \slidesources{\OOPSUnitVSources}
\end{frame}

\begin{frame}{Recall Before You Start}
\begin{itemize}[<+->]
  \item \textbf{Collection interfaces}: Utilities and wrappers only make sense once List, Set, and Map roles are already clear. Main reference: \texttt{Unit 05 / Lecture 01 slides/notes}.
  \item \textbf{Primitive and object types}: Wrapper classes connect directly to the type-conversion ideas from Unit 01. Main reference: \texttt{Unit 01 / Lecture 03 slides/notes}.
\end{itemize}
  \slidesources{\OOPSUnitVSources}
\end{frame}


\begin{frame}{Learning Outcomes}
  \begin{itemize}[<+->]
    \item Use common \texttt{Collections} utilities (sort, min/max, shuffle, reverse).
    \item Use wrapped collections (unmodifiable/synchronized/checked) appropriately.
    \item Avoid common wrapper/autoboxing pitfalls (and know class loading basics).
  \end{itemize}
  \slidesources{\OOPSUnitVSources}
\end{frame}

\begin{frame}{Key Terms to Know}
\small
\begin{columns}[t]
\column{0.48\textwidth}
\begin{itemize}
  \item \textbf{Utility class}: class with static helper methods (e.g., \texttt{Collections}).
  \item \textbf{Unmodifiable view}: wrapper that blocks structural modification.
  \item \textbf{Synchronized wrapper}: wrapper that adds basic synchronization around operations.
  \item \textbf{Checked wrapper}: wrapper that enforces runtime type checks on adds.
\end{itemize}
\column{0.48\textwidth}
\begin{itemize}
  \item \textbf{Wrapper class}: object form of a primitive (e.g., \texttt{Integer} for \texttt{int}).
  \item \textbf{Autoboxing/unboxing}: automatic conversion between primitive and wrapper.
  \item \textbf{Class loading}: JVM loads class bytecode into memory when needed.
\end{itemize}
\end{columns}
  \slidesources{\OOPSUnitVSources}
\end{frame}

\begin{frame}{Study Flow for This Lecture}
\begin{enumerate}[<+->]
  \item Refresh the prerequisite bridge before reading new ideas in isolation.
  \item Study the main build in this order: \textit{\texttt{Collections} utilities (common ones)} $\rightarrow$ \textit{Wrapped collections (why/when)} $\rightarrow$ \textit{Wrapper classes + autoboxing pitfalls} $\rightarrow$ \textit{Class loading basics (intro)}.
  \item Trace the worked example and connect each line back to one concept frame.
  \item Attempt the mini exercises before moving to the recap and exit check.
\end{enumerate}
  \slidesources{\OOPSUnitVSources}
\end{frame}


\section{Key Topics}
\label{sec:concepts}

\begin{frame}{Unit 05: Collections and Wrapper Class}
  \begin{itemize}[<+->]
    \item Lecture focus: Wrapped Collections + `Collections` Utilities + Wrappers
  \end{itemize}
  \slidesources{\OOPSUnitVSources}
\end{frame}

\begin{frame}{Today: Roadmap}
  \begin{itemize}[<+->]
    \item `Collections` utilities: sort, min/max, shuffle (overview)
    \item Wrapped collections: unmodifiable/synchronized/checked (why/when)
    \item Wrapper classes + autoboxing pitfalls; class loading basics (syllabus)
  \end{itemize}
  \slidesources{\OOPSUnitVSources}
\end{frame}

\begin{frame}{\texttt{Collections} utilities (common)}
  \begin{itemize}[<+->]
    \item \texttt{sort}, \texttt{reverse}, \texttt{shuffle}
    \item \texttt{min}, \texttt{max}, \texttt{frequency}
    \item \texttt{binarySearch} (requires sorted list)
  \end{itemize}
  \slidesources{\OOPSUnitVSources}
\end{frame}

\begin{frame}{Wrapped collections (why)}
  \begin{itemize}[<+->]
    \item \texttt{unmodifiableList}: expose read-only view to callers.
    \item \texttt{synchronizedList}: basic thread-safe wrapper (older style).
    \item \texttt{checkedList}: runtime type checking for legacy/raw code.
  \end{itemize}
  \slidesources{\OOPSUnitVSources}
\end{frame}

\begin{frame}{Wrapper classes + autoboxing pitfalls}
  \begin{itemize}[<+->]
    \item Wrapper classes: \texttt{Integer}, \texttt{Double}, ...
    \item Needed because generics require reference types.
    \item Pitfall: unboxing \texttt{null} $\rightarrow$ \texttt{NullPointerException}.
    \item Pitfall: prefer \texttt{equals()} over \texttt{==} for wrappers/strings.
  \end{itemize}
  \slidesources{\OOPSUnitVSources}
\end{frame}

\begin{frame}{Checkpoint and Reflection}
\begin{block}{Reflection prompt}
Where would \textit{Wrapped Collections + `Collections` Utilities + Wrappers} matter most in Campus Resource Manager, and which design choice would you justify using this lecture's ideas?
\end{block}
\begin{block}{Quick self-check}
Which part needs one more example before you move on: \textit{\texttt{Collections} utilities (common ones)}, \textit{Wrapped collections (why/when)}, the worked example, or the recap?
\end{block}
  \slidesources{\OOPSUnitVSources}
\end{frame}

\section{Demo / Example}
\label{sec:demo}

\begin{frame}[fragile,shrink=8]{Example: utilities + wrappers}
\begin{lstlisting}
import java.util.*;
public class Demo {
    public static void main(String[] args) {
        List<Integer> a = new ArrayList<>(Arrays.asList(3, 1, 2));
        Collections.sort(a);
        System.out.println("Sorted: " + a);

        List<Integer> ro = Collections.unmodifiableList(a);
        // ro.add(99); // UnsupportedOperationException

        Integer x = null;
        // int y = x; // NullPointerException (unboxing null)
    }
}
\end{lstlisting}
  \slidesources{\OOPSUnitVSources}
\end{frame}

\begin{frame}{Mini Exercises (2--3 mins)}
  \begin{enumerate}[<+->]
    \item What exception occurs when modifying an unmodifiable list?
    \item Why do we need \texttt{Integer} instead of \texttt{int} in lists?
    \item What happens when you unbox \texttt{Integer x = null}?
  \end{enumerate}
  \slidesources{\OOPSUnitVSources}
\end{frame}

\begin{frame}{Watch-outs}
\begin{itemize}[<+->]
  \item Assuming wrapper objects behave exactly like primitives in every context.
  \item Treating synchronized wrappers as a full concurrency design.
  \item Forgetting that an unmodifiable view can still reflect external changes to the backing collection.
\end{itemize}
  \slidesources{\OOPSUnitVSources}
\end{frame}

\section{Summary}
\label{sec:summary}

\begin{frame}{Key Takeaways}
  \begin{itemize}[<+->]
    \item \texttt{Collections} provides standard algorithms for collections.
    \item Wrappers enforce policies (read-only, sync, type checks).
    \item Wrapper classes enable generics; watch autoboxing pitfalls.
  \end{itemize}
  \slidesources{\OOPSUnitVSources}
\end{frame}

\begin{frame}{Checkpoint Questions}
  \begin{enumerate}[<+->]
    \item Difference: \texttt{Collection} vs \texttt{Collections}?
    \item When do you use an unmodifiable wrapper?
    \item Give one autoboxing pitfall.
  \end{enumerate}
  \slidesources{\OOPSUnitVSources}
\end{frame}

\begin{frame}{Next Study Steps}
\begin{itemize}[<+->]
  \item Revisit the recall references if any older idea still feels weak.
  \item Rework the lecture example without looking at the notes first.
  \item Attempt the self-study practice and then answer the exit questions in full sentences.
  \item Apply the idea once inside Campus Resource Manager to make the concept stick.
\end{itemize}
  \slidesources{\OOPSUnitVSources}
\end{frame}

\begin{frame}{Next Unit Preview}
  \textbf{Next:} Unit 06 -- Capstone Project
  \vspace{0.6em}
  \begin{itemize}[<+->]
    \item We will apply OOP + collections + Swing + JDBC in a complete project.
  \end{itemize}
  \slidesources{\OOPSUnitVSources}
\end{frame}

\end{document}
